\section{Arithmetic dynamics before the encoding}
\label{sec:arithmetic-dynamics}

The numen formalism does not begin with a transform.  It begins with an
ordinary discrete dynamical system and with the elementary distinction
between a bounded recurrent trajectory and a trajectory which escapes every
bounded subset.  This section fixes those notions before any coding is
introduced.

\subsection{Forward orbits and eventual repetition}

Let $Y$ be a set and let $T:Y\to Y$.  Write $T^{\circ 0}=\operatorname{id}_Y$
and $T^{\circ(n+1)}=T\circ T^{\circ n}$.

\begin{definition}
For $y\in Y$, the \emph{forward orbit} is
\[
  \mathcal O_T^+(y)=\bigl(T^{\circ n}(y)\bigr)_{n\geq 0}.
\]
The point $y$ is \emph{periodic} if $T^{\circ m}(y)=y$ for some $m\geq1$,
and \emph{preperiodic} if $T^{\circ a}(y)$ is periodic for some $a\geq0$.
It is \emph{strictly preperiodic} if it is preperiodic but not periodic.
The least positive period of a periodic point is its \emph{exact period}.
\end{definition}

\begin{lemma}[eventual repetition]
\label{lem:eventual-repetition}
For $y\in Y$, the following are equivalent.
\begin{enumerate}[label=\textup{(\roman*)}]
\item $y$ is preperiodic;
\item $T^{\circ m}(y)=T^{\circ n}(y)$ for some $0\leq m<n$;
\item the set $\{T^{\circ n}(y):n\geq0\}$ is finite.
\end{enumerate}
\end{lemma}

\begin{proof}
If $T^{\circ m}(y)=T^{\circ n}(y)$ with $m<n$, then
$T^{\circ m}(y)$ is periodic of period dividing $n-m$.  This proves
(ii)$\Rightarrow$(i).  A preperiodic orbit consists of a finite initial tail
followed by a finite cycle, proving (i)$\Rightarrow$(iii).  Finally, an
infinite sequence taking values in a finite set repeats a value, so
(iii)$\Rightarrow$(ii).
\end{proof}

The relation used in Siegel's latest manual is
\[
 x\sim_T y
 \quad\Longleftrightarrow\quad
 T^{\circ m}(x)=T^{\circ n}(y)
 \quad\text{for some }m,n\geq0.
\]
It is an equivalence relation: symmetry and reflexivity are immediate, and
transitivity follows by iterating the two common-forward-image equalities
until their middle iterates agree.  Its equivalence classes are the weak
connected components of the functional graph of $T$.

\subsection{The compact-bounded dichotomy}

Suppose now that $Y$ is contained in a normed finite-dimensional real vector
space $V$.  In the arithmetic applications the unconditional locally finite
model is the image of the ring of integers of a number field under its full
Minkowski embedding.  The image under one archimedean embedding need not be
locally finite when the field has degree greater than one.  A subset
$Y\subset V$ is \emph{locally finite} when
$Y\cap B$ is finite for every bounded $B\subset V$.

\begin{proposition}[the full Minkowski image is a lattice]
\label{prop:minkowski-lattice}
Let $K$ be a number field of degree $d=r_1+2r_2$.  Choose its real
embeddings $\sigma_1,\ldots,\sigma_{r_1}$ and one embedding
$\tau_j$ from each pair of complex-conjugate embeddings.  Then
\[
 \iota_K:K\longrightarrow\R^{r_1}\times\C^{r_2}\cong\R^d,
 \qquad
 x\longmapsto
 (\sigma_1(x),\ldots,\sigma_{r_1}(x),
  \tau_1(x),\ldots,\tau_{r_2}(x))
\]
maps $\OK$ onto a full lattice.  In particular, $\iota_K(\OK)$ is locally
finite.
\end{proposition}

\begin{proof}
Choose an integral basis $\omega_1,\ldots,\omega_d$ of $\OK$.  The matrix
whose columns are the images of these basis elements under all $d$
embeddings has nonzero determinant: its squared determinant is the nonzero
field discriminant of the basis.  Replacing each conjugate pair of complex
coordinates by its real and imaginary parts is an invertible real-linear
change of coordinates.  Hence the real-linear map
\[
 A:\R^d\longrightarrow\R^{r_1}\times\C^{r_2},
 \qquad e_j\longmapsto\iota_K(\omega_j),
\]
is an isomorphism, and
$\iota_K(\OK)=A(\Z^d)$.  If $B$ is bounded, then $A^{-1}(B)$ is bounded and
therefore meets $\Z^d$ in finitely many points.  Thus $B$ meets
$\iota_K(\OK)$ in finitely many points.
\end{proof}

\begin{proposition}[preperiodic or escaping]
\label{prop:preperiodic-or-escaping}
Let $Y\subset V$ be locally finite and let $T:Y\to Y$.  For every $y\in Y$
exactly one of the following occurs:
\begin{enumerate}[label=\textup{(\roman*)}]
\item $y$ is preperiodic;
\item $\|T^{\circ n}(y)\|\to\infty$ as $n\to\infty$.
\end{enumerate}
\end{proposition}

\begin{proof}
The alternatives are disjoint because a preperiodic orbit is finite.  If
$y$ is not preperiodic, \cref{lem:eventual-repetition} says that all members
of its forward orbit are distinct.  Fix $R>0$.  The set
$Y\cap\{v:\|v\|\leq R\}$ is finite, so only finitely many of those distinct
orbit points can lie in it.  Hence there is $N_R$ such that
$\|T^{\circ n}(y)\|>R$ for all $n\geq N_R$.  Since $R$ was arbitrary, the
norms tend to infinity.
\end{proof}

\begin{sourceissue}[the discreteness hypothesis in the 2026 statement]
The source states \cref{prop:preperiodic-or-escaping} for a uniformly
discrete subset of an arbitrary valued field and infers that every bounded
ball meets that subset in finitely many points.  Uniform separation alone
does not imply this in a non-proper metric space.  The proof is valid under
the locally finite hypothesis above; in particular it is valid for the full
Minkowski lattice.  It is not automatically valid after projection to one
archimedean embedding, where bounded sets may contain infinitely many
algebraic integers.  This edition does not silently retain either the more
general but insufficient source hypothesis or that single-embedding
specialization without an added discreteness assumption.
\end{sourceissue}

\subsection{The Collatz and shortened
\texorpdfstring{$qx+1$}{qx+1} maps}

The classical Collatz map on $\Z$ is
\[
 C(n)=
 \begin{cases}
 n/2,&n\equiv0\pmod2,\\
 3n+1,&n\equiv1\pmod2.
 \end{cases}
\]
Its odd branch is followed by an even value.  Compressing that forced even
step gives the shortened map
\[
 T_3(n)=
 \begin{cases}
 n/2,&n\equiv0\pmod2,\\
 (3n+1)/2,&n\equiv1\pmod2.
 \end{cases}
\]
More generally, for an odd integer $q$, put
\begin{equation}
\label{eq:Tq}
 T_q(n)=
 \begin{cases}
 n/2,&n\equiv0\pmod2,\\
 (qn+1)/2,&n\equiv1\pmod2.
 \end{cases}
\end{equation}
The restriction to odd $q$ is structural: it makes the second expression an
integer exactly on the odd residue class.  The parity itinerary of $n$ is
\[
 \varepsilon_k(n)=[T_q^{\circ k}(n)]_2\in\{0,1\},\qquad k\geq0,
\]
where $[a]_2$ denotes the standard residue representative.  For a finite
orbit segment of length $m$ we package the itinerary as
\[
 \boldsymbol\varepsilon_m(n)
   =(\varepsilon_0(n),\ldots,\varepsilon_{m-1}(n)).
\]

There are two orders in the corpus which must not be conflated.  The entries
of an itinerary are chronological: the branch indexed by $\varepsilon_0$ is
applied first.  The numen word $\mathbf j=(j_0,\ldots,j_{m-1})$ is used in
the ordered affine composition
\[
 H_{\mathbf j}=H_{j_0}\circ\cdots\circ H_{j_{m-1}},
\]
so the rightmost branch acts first.  Consequently the numen word attached to
a chronological orbit segment is its reversal unless a statement explicitly
adopts the opposite composition convention.  Every order-sensitive formula
below is derived rather than inferred from prose.

\subsection{What the coding is required to preserve}

The goal is not merely to replace an integer orbit by a digit string.  A
usable coding must simultaneously retain:
\begin{enumerate}[label=\textup{(\alph*)}]
\item the residue class which makes each branch the genuine branch of the
      piecewise map;
\item the linear part of a branch word;
\item the translation part of that branch word;
\item convergence at a specified completion when an infinite digit string
      is substituted; and
\item enough branch correctness to turn a fixed point of a formal branch
      composition back into a periodic point of the original map.
\end{enumerate}
Hydra maps provide (a), the pair $(M_H,X_H)$ provides (b)--(c), the extension
lemma provides (d), and integrality supplies the principal mechanism for (e).
